\chapter{HAI Robotics K50 Reference}

\section*{ii. Document 2}

The proposed AMR system is inspired by the K50 Fast-Transit Companion AMR developed by HAI Robotics. The robot is designed to operate as a fast-transit mobile platform capable of interacting directly with standardized warehouse racks and handling totes through a structured backpack mechanism.

\section*{Quick and convenient AMR features}

\begin{itemize}
    \item Maximum speed of 4 m/s.
    \item Maximum jacking height of 0.75 m suitable for ergonomic workstations.
    \item LiDAR remote scanning.
\end{itemize}

\section*{Mechanical and motion specifications}

\begin{table}[H]
\centering
\begin{tabularx}{\textwidth}{>{\bfseries}L{0.32\textwidth}Y}
\toprule
Parameter & Value \\
\midrule
Dimensions (L $\times$ W $\times$ H) & 658 $\times$ 460 $\times$ 330 mm \\
Net weight & 81 kg \\
Rated load & 30/50 kg \\
Case type mode & Carton, tote, and other materials with regular shapes \\
Driving mode & Differential-wheeled \\
Navigation & DM code + IMU \\
Driving speed & $\leq$ 4.0 m/s \\
Stop accuracy & $\pm$10 mm \\
Travel direction & Forward and backward, rotation in place, and turning \\
\bottomrule
\end{tabularx}
\end{table}

\section*{Functional description}

The core operational concept is based on a three-step rack interaction workflow: precise positioning under rack, tote securing through a backpack mode mechanism, and high-speed transport and delivery to workstation.

\subsection*{Precise positioning under rack}

The AMR navigates autonomously within warehouse aisles and performs precise alignment beneath storage racks equipped with low-height guide rails, approximately 0.43 m from the ground. This phase requires high-precision localization, fine motion control near racks, and accurate rack alignment detection.

\subsection*{Tote securing}

After correct positioning, the AMR engages a mechanical lifting and securing mechanism composed of twin cradling supports or fork-like arms. The mechanism engages the tote from the underside, stabilizes the load laterally, locks the tote securely onto the robot chassis, and transforms the tote into a backpack-style carried load.

\subsection*{Transport and delivery}

Once the tote is secured, the AMR transitions into fast-transit mode, transporting the loaded tote to a designated workstation or drop-off zone. Upon arrival, the robot navigates to a predefined delivery point, raises or positions the tote at an ergonomic height, and allows efficient manual picking or unloading operations.
